\section{Scope, Edition, and Qualifications}

\subsection{What this work is}

A reference to the \emph{norms} of the postconciliar liturgical year as they are read for the
dioceses of the United States of America under the Roman Missal, Third Edition. Its corpus is the
governing texts and the temporal cycle: the structure of the year, the grades of celebration, the
table of liturgical days, occurrence and transfer, the movable anchors, the numbering of Ordinary
Time and its omitted week, the Sunday and weekday Lectionary cycles, and the law by which national,
diocesan and church calendars are composed from the general one.

It is a study reference. It is not an \latin{Ordo}, not a liturgical book, not an act of any
authority, and not a substitute for the competent annual calendar. Where a computation in this
reference disagrees with the competent calendar of a diocese or of a conference of bishops, the
competent calendar governs and the disagreement should be reported, not resolved by preference.

\subsection{What it deliberately does not contain}

\begin{itemize}
\item The fixed universal inventory --- the dated entries of the General Roman Calendar --- and the
recurring national overlay for the United States. Individual entries appear only where a rule is
being illustrated. No completeness is claimed for any list of celebrations in this document, and the
term ``complete calendar'' is nowhere used of it.
\item Any diocesan, parish, cathedral, titular, dedication, patronal or religious-family calendar.
\item Any prayer, reading, antiphon, chant, preface or formulary. The rights position below explains
why.
\item A resolved calendar for any year. Section~\ref{sec:years} works three years to verify rules; a
worked year is evidence for a rule, not an \latin{Ordo} for use.
\item The Liturgy of the Hours beyond the points at which the Universal Norms themselves extend a
calendar rule to it. The \latin{Institutio generalis de Liturgia Horarum} was not read for this
reference, and Section~\ref{sec:synthesis} records that as a limit on one of its arguments.
\item The history of the reform. The dated timeline in Section~\ref{sec:timeline} is a list of governing acts, not a
narrative of how the calendar came to be.
\end{itemize}

\subsection{Editions and witnesses actually used}

\begin{ruletable}
\textbf{\latin{Normae universales de anno liturgico et de calendario}} & Latin, as printed in
\latin{Missale Romanum}, \latin{editio typica tertia} (2002). Read complete, nn.~1--61 with the table
of liturgical days.\\
\textbf{\latin{Institutio generalis Missalis Romani}} & Latin, in the same 2002 edition. Read at the
articles cited: 53, 68, 353--355, 357--358, 363, 372, 374--378.\\
\textbf{\latin{Missale Romanum}, 2002} & Read at the rubrics of Ordinary Time, the formularies
heading weeks I and XXXIV, the rubric at the Baptism of the Lord and its parallel in the General
Roman Calendar, the heading of the weekday Masses of Christmas Time, the rubric at the Epiphany
concerning the announcement of movable feasts, the movable entries appended to December and June in
the General Roman Calendar, and the \latin{Tabella temporaria}.\\
\textbf{\latin{Ordo lectionum Missae}} & Latin, \latin{editio typica altera} (1981). Read at
\latin{Praenotanda} 65, 66 with footnote 102, 69, 103, 104 with footnotes 115 and 116, and at the
decrees of 25 May 1969 and 21 January 1981 in the front matter.\\
\textbf{\emph{Sacrosanctum Concilium}} & Latin. Articles 102--111 and the appended declaration on the
revision of the calendar.\\
\textbf{\latin{Codex Iuris Canonici} (1983)} & Latin, canon 1246, in the Holy See's web presentation
of Book~IV.\\
\textbf{Instruction \emph{Calendaria particularia}} & Latin, Acta Apostolicae Sedis 62 (1970)
651--663. Read at Chapter~I nn.~1--6, Chapter~II nn.~7--27, and Chapter~III n.~36.\\
\textbf{\emph{Liturgical Calendar for the Dioceses of the United States of America}} & Editions 2026,
2027 and 2028, published by the Secretariat of Divine Worship of the United States Conference of
Catholic Bishops. Read at the front-matter tables and notes and at the daily entries needed to fix
the Ordinary Time week numbering.\\
\end{ruletable}

\subsection{Language, orthography and rendering}

Latin is quoted as it stands in the witness used, including its orthography: the 2002 Missal prints
the \ae{} and \oe{} ligatures where the 1981 Order of Readings prints the letters separately, and
this reference preserves each. Roman numerals for weeks of Ordinary Time follow the Missal's own
usage.

Every English rendering of a Latin norm in this document is a project working rendering, labelled as
such at the point of use. None is an approved translation, none may be used liturgically, and none
should be cited as the text of the norm. Where an approved English translation exists --- the
Universal Norms and the Introduction to the Lectionary are published in approved English by several
conferences --- this reference deliberately did not quote it, for the rights reason below, and worked
from the Latin instead.

Terms are used in the postconciliar sense throughout. \emph{Solemnity}, \emph{feast},
\emph{memorial} and \emph{optional memorial} are the four grades of \loc{nualc}~10 and 14; no term of
the 1960/1962 system is a synonym for any of them, and no equivalence between the two systems is
asserted anywhere in this work.

\subsection{Method and source hierarchy}

The order of authority applied throughout is: the promulgating act; the Latin typical edition of the
book in which the norm stands; the applicable universal instruction; the approved territorial
calendar; and, last, computation. Computation is never treated as its own witness. Every dated
result published here was computed from the norms and then checked against a dated official witness,
and the checks are reported in Section~\ref{sec:years} whether or not they succeeded.

Claims are kept in distinct states. Quoted Latin is verified source text read in the identified
artifact at the recorded locus. Statements introduced as what a source ``says'' at a named number are
checked paraphrase. Statements about how provisions combine --- for instance that the omission costs
the Lectionary ten weekday sets in a thirty-three-week year --- are source-grounded synthesis, and
each rests on named loci. The two blocks in Section~\ref{sec:synthesis} are labelled project
judgments and are the only places where this reference argues beyond its sources; each states its
strongest counterargument and what would change it.

\subsection{Material uncertainties and honest ceilings}

\begin{enumerate}
\item \textbf{No licensed altar book was collated.} The Latin Missal text used here was read in a
digitally typeset secondary reproduction identified in the source records, not in a licensed printed
\latin{editio typica} and not in a page-image facsimile. That artifact is demonstrably incomplete:
the contents of its first appendix are absent, and the \latin{Tabella temporaria} has lost its column
structure. Page-level identity with the printed edition is not claimed anywhere in this work.
\item \textbf{The leap-year defect is unresolved.} Section~\ref{sec:tabella} reports a systematic
one-day error in the February dates of five leap-year rows of the \latin{Tabella temporaria} as read
in that artifact. Whether it stands in the printed typical edition, was corrected in the 2008 emended
reprint, or was introduced in the reproduction was not determined.
\item \textbf{The 2008 emended reprint and the 2011 English edition were not examined.} Their
identity is reported from the Secretariat of Divine Worship's calendar. Any claim in this reference
about what the 2011 English altar book prints would be unsupported, and none is made.
\item \textbf{The 1969 first typical edition of the Order of Readings was not examined.} Its own
promulgating decree of 21 January 1981 records that the \latin{Praenotanda} were enlarged for the
second edition, so no claim is made here that any cited number or footnote stood in 1969.
\item \textbf{The \latin{Institutio generalis de Liturgia Horarum} was not read.} The Liturgy of the
Hours is treated only where \loc{nualc}~44 itself extends a rule to it, and the gap is named in
Section~\ref{sec:synthesis}.
\item \textbf{No ecclesiastical statement of the Easter computus was located.} The 2002 typical
edition contains none; the conciliar declaration presupposes the Gregorian computation without
stating it. The rule is stated in Section~\ref{sec:movable} as an inherited rule of the Gregorian calendar, and the
absence of a liturgical-book locus is a bounded negative result of the searches actually run, not a
proof that no such text exists.
\item \textbf{Two Roman acts are reported rather than read.} The 1998 notification on the Immaculate
Heart, protocol 2671/98/L, and the decrees inserting post-2002 additions into the General Roman
Calendar are known here only through the Secretariat's calendar. Their own texts were not reached at
the Holy See's portal on 25 July 2026 at the paths tried.
\item \textbf{The United States complementary norm to canon 1246~\S2 was not read at its own source.}
Its effect is reported from the Secretariat's calendar for 2026. The Conference's page carrying the
norm returned an access refusal to this session on 25 July 2026.
\item \textbf{No claim of completeness for any inventory.} No amendment to the table of liturgical
days after 1969 was located; that is a bounded negative result over the sources actually consulted.
The recurring universal and national inventories of celebrations are outside this leaf's corpus
altogether.
\item \textbf{Territorial state is as of the cutoff.} The provincial distribution of the Thursday and
Sunday Ascension in the United States is reported from the 2026, 2027 and 2028 calendars and can
change by decision of the competent authority without notice to this document.
\end{enumerate}

\subsection{Rights}

No prayer, reading, antiphon, chant, preface or formulary is reproduced in this work. The English
translation of the Roman Missal for the dioceses of the United States and the English \emph{Lectionary
for Mass} are under copyright and are not quoted here at any length or paraphrased as substitutes;
where their content is relevant it is described and cited by locus. The annual calendars of the
Secretariat of Divine Worship carry an express reservation of rights and are used here only as
factual witnesses to dates, week numbers, cycle labels and territorial decisions, with citation.

The Latin normative texts quoted --- the Universal Norms, the General Instruction, the rubrics of the
Missal, the \latin{Praenotanda} and decrees of the Order of Readings, \emph{Sacrosanctum Concilium},
the Code, and the 1970 instruction --- are official acts of the Church quoted for study at the loci
named. Quotation here asserts no permission beyond that use and relicenses nothing. The exact
artifacts through which they were read are identified in the repository's source records with their
rights status; none of their bytes is redistributed by this publication.

\subsection{Review state}

Internal source and computation review is complete: every Latin quotation was read at the recorded
locus in the identified artifact, the passages from the 1981 Order of Readings were additionally
collated against rendered page images, and every dated result was recomputed and checked against an
official witness where one exists. That is internal checking, not independent review. Independent
liturgical, canonical, territorial-law, Latin-philological and specialist review remain outstanding,
as does work-specific rights review. This document carries no ecclesiastical approval of any kind and
is not offered as one.
